\section{Introduction}

\subsection{Motivation}

A memorable viewing experience includes the people in the room. Friends may laugh at the same unexpected line, a couple may gasp at a plot twist, or a home party may erupt during a tense scene. These reactions can carry the mood of the occasion more vividly than a photograph taken afterward. They also happen without warning. Reaching for a phone after the moment has passed cannot recover the initial reaction, while recording an entire evening leaves users with material to search and edit.

We propose Encore Moment to help viewers preserve these spontaneous moments while remaining engaged with one another and the screen. Research on preserving everyday laughter describes both the enjoyment of revisiting it with loved ones and the tensions involved in collecting it \cite{ref1}. This supports treating reactions as personal memories, while leaving users in control of what is retained. It does not establish demand for a TV service; the viewing scenarios in this proposal are intended use cases to be validated with users.

\subsection{Target Users and Viewing Contexts}

Friends watching together and guests at a home party may want a short record of the group's response without assigning someone to act as a camera operator. Couples on a movie date may value a shared expression or quiet moment of surprise, making unobtrusive operation and clear recording choices especially important. Families may wish to revisit reactions across ordinary evenings as well as special occasions. Across these settings, the value lies in who was present and what happened together, rather than simply accumulating recordings.

Streamers and reaction-content creators form an additional audience. They may need to find spontaneous reactions after a viewing session without reviewing the entire recording. Their editing needs differ from private social viewing, but both groups benefit from useful event markers and reviewable clips. Encore Moment focuses on finding and preserving candidate moments; users decide whether a saved moment is suitable to keep or use elsewhere.

\subsection{User and Client Needs}

Viewers need relevant highlights, enough context to remember the scene, and controls that are easy to understand during viewing. A film's laugh track or scream should not be mistaken for a person in the room. Users should know when sensing is active, be able to stop it, and review or delete saved clips. Reaction intensity can help prioritize review, but it cannot determine a person's emotional state or the personal importance of a memory.

For this proposal, the assumed client is a smart TV service provider seeking a useful viewing feature that fits an existing TV experience. LG smart TV integration is the intended project context. The client needs a clear feature scope, understandable privacy controls, and honest statements about supported devices and content access. These are proposed client needs, not confirmed contractual requirements or completed interview findings.
